\documentclass[12pt,a4paper]{article}
% Preámbulo del informe de la Etapa 3: el de la Etapa 2 sin la notación de la
% gramática. Computer Modern a 12pt, como los documentos de la cátedra, y todo en
% negro.
\usepackage[T1]{fontenc}
\usepackage[utf8]{inputenc}
\usepackage{lmodern}
\usepackage{amsmath,amssymb}
\usepackage[spanish,es-noquoting,es-noshorthands,es-tabla]{babel}
\usepackage[a4paper,margin=2.6cm,headsep=14pt]{geometry}
\usepackage{microtype}
\usepackage{booktabs}
\usepackage{longtable}
\usepackage{array}
\usepackage{xcolor}
\usepackage{fancyhdr}
\usepackage{titlesec}
\usepackage{enumitem}
\usepackage{fancyvrb}
\usepackage[hidelinks,bookmarks=true]{hyperref}

\usepackage{float}
\usepackage{tikz}
\usetikzlibrary{arrows.meta}
\usepackage{pgf-umlcd}
% Generalización e implementación de UML: triángulo abierto, línea llena o de trazos.
\tikzset{generaliza/.style={draw=black, -{Triangle[open, length=3mm, width=3mm]}},
        implementa/.style={generaliza, dashed}}
\DeclareUnicodeCharacter{2014}{---}
\DeclareUnicodeCharacter{2026}{\ldots}

% --- Paleta ------------------------------------------------------------------
\definecolor{tinta}{rgb}{0,0,0}
\renewcommand{\umltextcolor}{black}
\renewcommand{\umlfillcolor}{white}
\renewcommand{\umldrawcolor}{black}

\color{tinta}

% --- Títulos ------------------------------------------------------------------
\titleformat{\section}{\Large\bfseries\color{tinta}}{\thesection}{0.7em}{}
\titleformat{\subsection}{\large\bfseries\color{tinta}}{\thesubsection}{0.7em}{}
\titleformat{\subsubsection}{\bfseries\color{tinta}}{\thesubsubsection}{0.7em}{}
\titlespacing*{\section}{0pt}{2.4ex plus 1ex minus .2ex}{1.3ex plus .2ex}

% --- Encabezado y pie ---------------------------------------------------------
\setlength{\headheight}{14pt}
\pagestyle{fancy}
\fancyhf{}
\renewcommand{\headrulewidth}{0.4pt}
\renewcommand{\headrule}{\hbox to\headwidth{\color{tinta}\leaders\hrule height \headrulewidth\hfill}}
\fancyhead[L]{\footnotesize\color{tinta}Compiladores e Intérpretes — Etapa 3}
\fancyhead[R]{\footnotesize\color{tinta}\nouppercase{\leftmark}}
\fancyfoot[C]{\footnotesize\color{tinta}\thepage}

% --- Bloques de consola --------------------------------------------------------
% \footnotesize y no \small: la línea más larga de una sesión de consola tiene
% 82 caracteres y a \small no entra en la caja de texto.
\DefineVerbatimEnvironment{consola}{Verbatim}%
  {fontsize=\footnotesize,xleftmargin=1.1em,formatcom=\color{tinta}}

\setlength{\emergencystretch}{4.5em}
\tolerance=3000

\setlist[itemize]{leftmargin=1.3em,itemsep=0.25ex,topsep=0.5ex}
\setlist[enumerate]{leftmargin=1.5em,itemsep=0.25ex,topsep=0.5ex}
\setlist[description]{leftmargin=1.3em,itemsep=0.35ex,topsep=0.6ex}
\renewcommand{\arraystretch}{1.25}

\fancyhead[L]{\footnotesize\color{tinta}Compiladores e Intérpretes --- Etapa 5}
\fancyhead[R]{\footnotesize\color{tinta}Dudas}

\newcommand{\cod}[1]{\texttt{#1}}
\newcommand{\pregunta}[1]{\par\medskip\needspace{4\baselineskip}\noindent\textbf{#1}\par\smallskip}
\DefineVerbatimEnvironment{fuente}{Verbatim}%
  {fontsize=\small,xleftmargin=1.2em,frame=leftline,framerule=1.2pt,framesep=0.8em,numbers=left,numbersep=4pt}
\usepackage{needspace}
\titleformat{\section}{\Large\bfseries\raggedright}{\thesection}{0.7em}{}
\setlength{\parindent}{0pt}
\setlength{\parskip}{0.55ex}

\begin{document}
\thispagestyle{fancy}
\raggedright

{\LARGE\bfseries Consultas --- Etapa 5 (generación de código)\par}
\smallskip
{\color{tinta}Agustín Paul --- L.U. 147524\par}
\smallskip
{\color{tinta}\rule{\textwidth}{0.6pt}}
\medskip

Te paso las dudas de la Etapa 5. En cada una hay un ejemplo y lo que pienso hacer; con un «sí» o
con la otra opción me alcanza. «El documento» es \emph{Semántica de MiniJava -- Documento
completo} (v3.0) y «el apunte», el \emph{Apunte CeIVM 2026}. Las líneas arrancan en 1.

\pregunta{1. ¿Qué error doy si no hay \cod{main}, o si hay más de uno?}

El documento (§8 y §14) deja para esta etapa el control del punto de entrada, y el apunte (§15)
dice que todo programa declara un \cod{static void main()}. No encontré qué reportar:

\begin{fuente}
class Aa {
    void main() { }
}
\end{fuente}

Mi idea: cuenta sólo un \cod{main} estático, \cod{void} y sin parámetros, declarado en la clase (no
heredado). Si no hay ninguno, un error semántico en el fin de archivo, como hace el sintáctico
cuando falta algo al final: acá, \cod{[Error:\$|4]}, porque el archivo termina con un salto de
línea. Si hay dos, en clases distintas, el error va
en el nombre del segundo, como un nombre repetido.\\
Alternativas: reportar en otro token, o no controlarlo (en Java compila y falla recién al ejecutar).

\pregunta{2. Los errores de ejecución que controla mi código, ¿con qué mensaje terminan?}

\begin{fuente}
class Init {
    static void main() {
        var a = new int[2];
        a[5] = 1;
    }
}
\end{fuente}

El apunte (§12 y §17) pide controlar el índice fuera de rango y el tamaño negativo, y sugiere una
rutina que imprima un mensaje y haga \cod{HALT}. Pero tras un \cod{HALT} la CeIVM imprime «La
ejecución del programa finalizó exitosamente.», como en un programa sin errores. Mi idea: imprimir
\cod{Error en tiempo de ejecucion: indice fuera de rango} y terminar con \cod{HALT}.\\
¿Está bien así, o preferís que termine provocando un error de la propia VM, para que no diga
«exitosamente»? ¿Hay un formato de mensaje que esperen?

\pregunta{3. ¿Qué valor tienen los escapes en los literales?}

\begin{fuente}
class Init {
    static void main() {
        System.printCln('\n');
        System.printCln('\q');
        System.printSln("a\tb");
    }
}
\end{fuente}

Los aspectos léxicos (§2.3.2) aceptan \cod{'\textbackslash x'} con cualquier \cod{x}, pero no dicen
su valor, y para los Strings dicen que «el valor del literal corresponde a la cadena de caracteres
entre las comillas». Mi idea: como en Java, \cod{\textbackslash n} es el salto de línea,
\cod{\textbackslash t} el tabulador, etcétera, también dentro de los Strings; y como Java rechaza
\cod{'\textbackslash q'}, para los escapes que Java no tiene, el valor es el carácter mismo
(\cod{q}).\\
Alternativa: los Strings se toman crudos, y \cod{"a\textbackslash tb"} tiene cuatro caracteres.

\pregunta{4. Para confirmar: en una asignación a un atributo o a un arreglo, ¿evalúo primero el lado derecho?}

\begin{fuente}
class Init {
    int x;
    static Init obj() { debugPrint(1); return new Init(); }
    static int val() { debugPrint(2); return 5; }
    static void main() {
        obj().x = val();
    }
}
\end{fuente}

Sigo el esquema de la Clase 16 (slides 18 a 29) y del apunte (§7): primero el valor, después la
referencia, \cod{SWAP} y \cod{STOREREF}. Imprime 2 y después 1. En Java imprime 1 y después 2: el
lado izquierdo se evalúa primero.\\
Alternativa: hacerlo como Java.

\pregunta{5. Para confirmar: dos literales String iguales, ¿son el mismo objeto?}

\begin{fuente}
class Init {
    static void main() {
        System.printBln("hola" == "hola");
    }
}
\end{fuente}

El apunte (§13) aloja cada literal en el heap y el documento (§11.2) compara las referencias por
identidad, así que, creando el String cada vez que se evalúa el literal, imprime \cod{false}. En
Java imprime \cod{true}, porque los literales iguales son el mismo objeto.\\
Alternativa: crear cada literal distinto una sola vez, al arrancar, para que dé \cod{true} como en
Java.

\end{document}
